\documentclass[type=report,theme=light,language=de]{semio}
\title{Abbreviations Verify}
\author{Semio}
\date{\today}

\GlossaryDefine{Entwurfsfähigkeit}{Bauteil mit belastbaren Informationen für frühe Entwurfsentscheidungen.}
\GlossaryDefine{Semio}{Einheitliches Betriebssystem für gebaute Umgebung und Designprozesse.}

\AbbreviationDefine{KI}{Künstliche Intelligenz}{Oberbegriff für Computerprogramme, die Aufgaben lösen, für die normalerweise menschliches Denken nötig ist.}
\AbbreviationDefine{LLM}{Large-Language-Model}{Ein Computerprogramm, das mit sehr großen Textmengen trainiert wurde.}
\AbbreviationDefine{AP}{Arbeitspaket}{Ein klar abgegrenzter Teilbereich des Projekts.}

\begin{document}

\maketitle

\chapter{Ergebnisse}
Die \Gls{Entwurfsfähigkeit} wurde mithilfe von \Abbr{KI} untersucht.
\Gls{Semio} nutzt dabei ein \Abbr{LLM}.
Auf dieser Seite erscheint \Gls{Entwurfsfähigkeit} erneut und \Abbr{KI} noch einmal.

\chapter{Projektstand}
Das \Abbr{AP} fasst den aktuellen Stand zusammen.
\Gls{Semio} und \Abbr{LLM} werden hier in einem anderen Kapitel referenziert.

\appendix
\chapter{Verzeichnisse}
\listofglossaries
\listofabbreviations

\end{document}
